\documentclass[uplatex,dvipdfmx,a4paper,11pt]{jsarticle}
\input{../../../templates/note.tex}
\noteprint{数理物理1　第9章「離散フーリエ変換と周期境界の連成振動」}{2026-11-25}
%% note-moto	problems\数理物理1\shou09\数理物理1-9.tex
\begin{document}
\begin{kirokuichiran}
  \ichirangyou{1}{L1・計算}{◎◎◎○}{A}{2026-10-14 2回目 (4)記述}
  \ichirangyou{2}{L1・計算}{◎◎◎}{S}{2026-10-03}
  \ichirangyou{3}{L2・導出}{○▲△△}{C}{2026-10-14 2回目 (1)記述 (2)立式 (3)計算 (4)計算}
  \ichirangyou{4}{L2・計算}{○◎○◎}{B}{2026-10-12 2回目 ヒント1 (1)記述 (3)記述}
  \ichirangyou{\hyperlink{bt-m-097-ringinit}{5}}{L2・計算}{◎◎××}{C}{2026-10-03 (3)立式 (4)未着手}
  \ichirangyou{6}{L2・計算}{}{未着手}{}
  \ichirangyou{7}{L3・計算}{◎×△×××}{D}{2026-10-03 (2)未着手 (3)記述 (4)記述 (5)未着手 (6)未着手}
  \ichirangyou{8}{L3・導出}{}{未着手}{}
  \ichirangyou{9}{L3・計算}{}{未着手}{}
  \ichirangyou{10}{L3・計算}{}{未着手}{}
  \ichirangyou{11}{L4・計算}{}{未着手}{}
\end{kirokuichiran}

\youtenseiri
\begin{noteshikilist}
  \item 連成振動の運動方程式 % 出典: p.88
    \[ m\dv[2]{x_n}{t}=-\pdv{V}{x_n}=K(x_{n+1}-x_n)-K(x_n-x_{n-1}),\quad n=1,2,\dots,N \notetag{9.18} \]
\end{noteshikilist}

\mondaihen
\level{2}
% toubun: kiroku\toubun\数理物理1\097_2026-10-03.tex
\mondaibangou{5}
\begin{mondai}{097-ringinit}{計算}{輪状の連成振動の初期値問題と，$\omega=0$ のモード（重心運動）}
質量 $m$ の質点 4 個（$N=4$）がばね定数 $K$ のばねで輪状につながっている（$x_{n+4}=x_n$，運動方程式は \noteshiki{9.18}）。$U_{n\mu}=\ee^{\ii2\pi\mu n/4}/\sqrt{4}$ とする。時刻 $0$ に質点 1 だけ $x_1(0)=x_0$，他は $x_2(0)=x_3(0)=x_4(0)=0$ で，全質点の速度は 0 である。
  \begin{shoumon}
    \item $\omega_\mu$（$\mu=1,2,3,4$）を $K,m$ で表し，縮退しているモードと，$\omega=0$ のモードを示せ。
    \item 基準座標の初期値 $\hat{x}_\mu(0)$ と $\dv{\hat{x}_\mu}{t}(0)$ を求めよ。
    \item 次の $x_n(t)$ を計算し，$\omega_a=\sqrt{2K/m}$，$\omega_b=2\sqrt{K/m}$ を用いて $x_1,x_2,x_3,x_4$ を表せ。
      \[ x_n(t)=\sum_\mu U_{n\mu}\hat{x}_\mu(t) \]
    \item $t=0$ で初期条件を再現することと，次の $S(t)$ が一定であることを確かめ，$x_2(t)=x_4(t)$ となる理由を述べよ。
      \[ S(t)=\sum_nx_n(t) \]
  \end{shoumon}
\notekiroku{2026-10-03 C ◎◎××}
\begin{kotaeutsushi}{2026-10-03}
\toi{1}{◎}{}
（図：4個の質点を輪状につないだ図。質点1の横に $x_0$）
\[ m\dv[2]{x_n}{t}=K(x_{n+1}-x_n)-K(x_n-x_{n-1}) \]
\begin{align*}
  \omega_\mu&=\sqrt{\frac{2K\bigl(1-\cos\frac{2\pi\mu}{N}\bigr)}{m}}\\
  &=\sqrt{\frac{2K}{m}\Bigl(1-\cos\Bigl(\frac{\pi\mu}{2}\Bigr)\Bigr)}
\end{align*}
\begin{align*}
  \omega_1&=\sqrt{\frac{2K}{m}}\\
  \omega_2&=2\sqrt{\frac{K}{m}}\\
  \omega_3&=\sqrt{\frac{2K}{m}}\\
  \omega_4&=0\quad\fumei
\end{align*}
（$\omega_1$ から $\omega_3$ を括弧でまとめている）
\toi{2}{◎}{}
\begin{align*}
  \hat{x}_\mu&=\sum U^{*}_{n\mu}x_n\\
  \hat{x}_\mu(0)&=\sum U^{*}_{n\mu}x_n(0)\\
  \hat{x}_\mu(0)&=\frac{\kakimachigai{\ee^{-\ii2\pi\mu n/4}}{\ee^{-\ii2\pi\mu/4}}}{2}x_0\\
  &=\frac{x_0}{2}\ee^{-\ii\mu\pi/2}
\end{align*}
\[ \dv{\hat{x}_\mu}{t}=0 \]
\toi{3}{×}{立式}
\[ x_n(t)=\sum U_{n\mu}\hat{x}_\mu(t) \]
\begin{align*}
  x_1(t)&=\sum_\mu\frac{\ee^{\ii2\pi\mu/4}}{2}\hat{x}_\mu(t)\\
  &=\frac{1}{2}\cdot\frac{x_0}{2}\ayamari{\ee^{\ii2\pi\mu/2}}{\sum_\mu\cos\omega_\mu t}
\end{align*}
\naze{位相は $\ee^{\ii\pi\mu/2}\cdot\ee^{-\ii\pi\mu/2}=1$ で打ち消し合う。また各モードの時間変化を入れていない：$\mu=1,2,3$ は $\hat{x}_\mu(t)=\hat{x}_\mu(0)\cos\omega_\mu t$，$\mu=4$ は $A+Bt$ で初速 0 なので定数。これを入れて $x_n=\frac{x_0}{4}\sum_\mu\ee^{\ii\pi\mu(n-1)/2}\cos\omega_\mu t$ とし，$x_1=\frac{x_0}{4}(2\cos\omega_at+\cos\omega_bt+1)$，$x_2=x_4=\frac{x_0}{4}(1-\cos\omega_bt)$，$x_3=\frac{x_0}{4}(-2\cos\omega_at+\cos\omega_bt+1)$ となる。}
（下に続けて）
\[ \hat{x}_\mu=\frac{x_0}{2}\ee^{-\ii\mu\pi/2},\qquad\hat{x}_0=\frac{x_0}{2}\ee \]
\naze{ここで止まっている。}
\toi{4}{×}{未着手}
\mikai
\end{kotaeutsushi}
\begin{kizuki}
\item 基準座標で解くときは，戻す前に各モードの時間変化を入れる。初速 0 のモードは $\hat{x}_\mu(0)\cos\omega_\mu t$，$\omega=0$ のモードは $\ee^{\pm\ii\omega t}$ が使えないので $A+Bt$（初速 0 なら定数）とし，$\mu$ の和は残したまま書く。
\item $U_{1\mu}\hat{x}_\mu(0)$ の位相は $\ee^{\ii\pi\mu/2}\ee^{-\ii\pi\mu/2}=1$ で消える。指数をまとめるときは符号を足し合わせて確かめ，最後に $t=0$ で $x_1=x_0$ を再現するかで検算する。
\item (4) は $\sum_nx_n$ で $\mu=4$ の項だけが残る（$\sum_n\ee^{\ii\pi\mu(n-1)/2}$ が $\mu=4$ で 4，他で 0）ことを使う。
\end{kizuki}
\begin{tadashii}
\houshin 次の式で基準座標の初期値を求め，各モードを解いて $U$ で戻す。$\omega=0$ のモードは初速が 0 なら定数。
\[ \hat{x}_\mu(0)=\sum_n(U^{\dagger})_{\mu n}x_n(0) \]
\keisan
\shoumonkai{1} $N=4$ では
\[ \omega_\mu^{2}=\frac{2K}{m}\Bigl(1-\cos\frac{\pi\mu}{2}\Bigr) \]
なので $\mu=1$：$2K/m$，$\mu=2$：$4K/m$，$\mu=3$：$2K/m$，$\mu=4$：$0$。$\omega_1=\omega_3=\omega_a=\sqrt{2K/m}$（縮退），$\omega_2=\omega_b=2\sqrt{K/m}$，$\omega_4=0$。$\mu=4=N$ のモードは全質点が同位相で動く重心運動で復元力がない。
\shoumonkai{2} $(U^{\dagger})_{\mu n}=\ee^{-\ii2\pi\mu n/4}/2$ なので $\hat{x}_\mu(0)=(U^{\dagger})_{\mu1}x_0=\frac{x_0}{2}\ee^{-\ii\pi\mu/2}$。$\mu=1$：$-\ii x_0/2$，$\mu=2$：$-x_0/2$，$\mu=3$：$+\ii x_0/2$，$\mu=4$：$x_0/2$。速度は全て 0 なので $\dv{\hat{x}_\mu}{t}(0)=0$。
\shoumonkai{3} $\mu=1,2,3$ は $\hat{x}_\mu(t)=\hat{x}_\mu(0)\cos\omega_\mu t$。$\mu=4$ は $\hat{x}_4=A+Bt$ で初速 0 より $B=0$，$\hat{x}_4(t)=x_0/2$（$\cos0=1$ としてまとめてよい）。よって
\begin{align*}
  x_n(t)&=\sum_{\mu=1}^{4}\frac{\ee^{\ii\pi\mu n/2}}{2}\cdot\frac{x_0}{2}\ee^{-\ii\pi\mu/2}\cos\omega_\mu t=\frac{x_0}{4}\sum_{\mu=1}^{4}\ee^{\ii\pi\mu(n-1)/2}\cos\omega_\mu t\\
  &=\frac{x_0}{4}\Bigl[2\cos\frac{\pi(n-1)}{2}\cos\omega_at+(-1)^{n-1}\cos\omega_bt+1\Bigr]
\end{align*}
（$\mu=1,3$ の和 $\ii^{\,n-1}+(-\ii)^{n-1}=2\cos\frac{\pi(n-1)}{2}$，$\mu=2$ は $(-1)^{n-1}$，$\mu=4$ は $1$）。成分は
\begin{align*}
  x_1&=\frac{x_0}{4}(2\cos\omega_at+\cos\omega_bt+1),\qquad x_3=\frac{x_0}{4}(-2\cos\omega_at+\cos\omega_bt+1)\\
  x_2&=x_4=\frac{x_0}{4}(1-\cos\omega_bt)
\end{align*}
\shoumonkai{4} $t=0$：$x_1=x_0$，$x_2=x_4=0$，$x_3=\frac{x_0}{4}(-2+1+1)=0$ で初期条件を再現する。和は
\[ \sum_nx_n=\frac{x_0}{4}\cdot4=x_0 \]
（$\cos\omega_a$ と $\cos\omega_b$ の項は和で消える）で一定。これは $\mu=4$ のモード（重心）が動かないことに対応する。質点 2 と 4 は質点 1 から輪に沿って等距離にある鏡像の位置で，運動方程式と初期条件が鏡映で不変なので $x_2(t)=x_4(t)$。
\kotae{\shoumonkai{1} $\omega_1=\omega_3=\sqrt{2K/m}$，$\omega_2=2\sqrt{K/m}$，$\omega_4=0$ \shoumonkai{2} $\hat{x}_\mu(0)=\frac{x_0}{2}\ee^{-\ii\pi\mu/2}$，速度の初期値は 0 \shoumonkai{3} $x_1=\frac{x_0}{4}(2\cos\omega_at+\cos\omega_bt+1)$，$x_2=x_4=\frac{x_0}{4}(1-\cos\omega_bt)$，$x_3=\frac{x_0}{4}(-2\cos\omega_at+\cos\omega_bt+1)$ \shoumonkai{4} 本文のとおり（$\sum x_n=x_0$）}
\kenzan 運動方程式から $t=0$ での加速度を求める。$m\dv[2]{x_1}{t}=K(x_2-2x_1+x_4)=-2Kx_0$，$m\dv[2]{x_2}{t}=K(x_1-2x_2+x_3)=Kx_0$。答えを2回微分して $t=0$ とすると
\begin{align*}
  \dv[2]{x_1}{t}&=-\frac{x_0}{4}(2\omega_a^{2}+\omega_b^{2})=-\frac{x_0}{4}\cdot\frac{8K}{m}=-\frac{2Kx_0}{m}\\
  \dv[2]{x_2}{t}&=\frac{x_0}{4}\omega_b^{2}=\frac{Kx_0}{m}
\end{align*}
で一致する。
\tsumazuki $\mu=4$（$\omega=0$）のモードは $\ee^{\pm\ii\omega t}$ が使えない。初速 0 なら定数，初速があれば $A+Bt$ で一様に動き続ける。$\hat{x}_\mu(0)$ の位相 $\ee^{-\ii2\pi\mu/N}$ を落とさない。
\end{tadashii}
\end{mondai}

\noteowari
\end{document}
